\documentclass[12pt, a4paper]{article}
\usepackage[utf8]{inputenc}
\usepackage[ngerman]{babel}
\usepackage{geometry}
\usepackage{graphicx}
\usepackage{hyperref}
\usepackage{enumitem}
\usepackage{tabularx}

\geometry{left=2.5cm, right=2.5cm, top=2.5cm, bottom=2.5cm}

\title{Pflichtenheft}
\author{Fachinformatiker Anwendungsentwicklung}
\date{\today}
\begin{document}

%--- Deckblatt ---
\begin{titlepage}
    \centering
    \vspace*{2cm}
    {\Huge \textbf{Pflichtenheft}}\\[1.5cm]
    {\Large \textbf{Projektname:}\\[0.3cm]
    \textbf{\glqq Bookii-gestütztes Lagersystem für private Haushalte\grqq}}\\[2cm]

    \begin{tabular}{>{\bfseries}l l}
        Auftraggeber: & GFN \\
        Projektleitung: & Fachinformatiker Anwendungsentwicklung \\
        Datum: & \today \\
    \end{tabular}
    \vfill
\end{titlepage}

\setcounter{page}{1}
\tableofcontents
\newpage

\section{Zielbestimmung}
Dieses Pflichtenheft beschreibt die detaillierten technischen und fachlichen Anforderungen zur Umsetzung des im Projektantrag definierten Projekts \glqq Bookii-gestütztes Lagersystem für private Haushalte\grqq. Ziel ist die Entwicklung eines funktionalen Prototyps, der die Lagerverwaltung in einem privaten Haushalt mittels Tangible User Interfaces (TUI) durch den Bookii-Stift unterstützt.

\subsection{Muss-Kriterien}
\begin{itemize}
    \item Generierung und Ausdruck von Bookii-kompatiblen Punktcodes.
    \item Zuweisung einer eindeutigen digitalen ID (Startcode) zu physischen Artikeln via Punktcode.
    \item Erfassung von Artikeln in einer Datenbank (Hinzufügen, Identifizieren).
    \item Systemarchitektur bestehend aus bedrucktem Papier (Frontend), Bookii-Stift, und Backend-Logik (Java).
    \item Persistente Speicherung der Lagerdaten in einer relationalen Datenbank (MariaDB).
\end{itemize}

\subsection{Soll-Kriterien}
\begin{itemize}
    \item Vorschlagssoftware für Lagerorte.
    \item Kurze Reaktionszeiten ($<$ 2 Sekunden) für eine flüssige Interaktion.
    \item Strukturierter und versionierter Quellcode (Git).
\end{itemize}

\subsection{Kann-Kriterien (optional)}
\begin{itemize}
    \item Erweiterte Statistiken über Lagerbestände.
    \item Weboberfläche zur visuellen Verwaltung parallel zum Stift.
\end{itemize}

\subsection{Abgrenzung}
Nicht Teil des Projekts ist die kommerzielle Veröffentlichung, das Entschlüsseln proprietärer Malbuchcodes oder die Neuentwicklung der Bookii-Stift-Hardware.

\section{Produkteinsatz}
\subsection{Anwendungsbereiche}
Private Haushalte und Demonstrationsumgebungen im betrieblichen Ausbildungskontext der GFN.
\subsection{Zielgruppen}
Endanwender, die ihr Inventar verwalten möchten, ohne durchgehend ein klassisches Display (Smartphone/PC) bedienen zu müssen.

\section{Produktübersicht / Systemarchitektur}
Das System besteht aus drei Hauptkomponenten:
\begin{enumerate}
    \item \textbf{Physisches Interface:} Bedrucktes Papier mit spezifischen Punktmustern (Little-Endian und Big-Endian Codierungen).
    \item \textbf{Eingabegerät:} Der Bookii-Stift, der die Codes optisch erfasst und verarbeitet.
    \item \textbf{Backend:} Eine Java-Anwendung, die mit einer MariaDB-Datenbank kommuniziert, um die Logik der Lagerverwaltung abzubilden.
\end{enumerate}

\section{Produktfunktionen}
\begin{itemize}
    \item \textbf{F01 - Code-Generierung:} Erzeugung von ausdruckbaren PNG-Bildern mit definierten Punktcodes in verschiedenen Formaten (einfach, zweifach, dreifach mit Bilderbuchcode).
    \item \textbf{F02 - Artikelregistrierung:} Ein neuer Artikel wird durch Antippen des generierten Codes mit dem Bookii-Stift im System registriert.
    \item \textbf{F03 - Artikelverwaltung:} Bearbeiten und Löschen von registrierten Artikeln.
    \item \textbf{F04 - Persistenz:} Speichern der Metadaten (ID, Name, Ort) in der Datenbank.
\end{itemize}

\section{Produktdaten}
Das System muss folgende Daten dauerhaft speichern (in MariaDB):
\begin{itemize}
    \item \textbf{Artikel:} ID (generiert aus dem Code), Name, Kategorie.
    \item \textbf{Lagerorte:} ID, Name, Beschreibung.
    \item \textbf{Lagerbestand:} Verknüpfung von Artikel-ID, Lagerort-ID und Anzahl.
\end{itemize}

\section{Produktleistungen}
\begin{itemize}
    \item \textbf{Zuverlässigkeit:} Das System muss eingehende Signale des Stiftes fehlerfrei den Datenbankeinträgen zuordnen.
    \item \textbf{Reaktionszeit:} Audio- oder Systemfeedback innerhalb von 2 Sekunden nach dem Scannen.
\end{itemize}

\section{Qualitätsanforderungen}
\begin{itemize}
    \item Modulare und erweiterbare Java-Architektur.
    \item Code-Dokumentation gemäß Best Practices.
    \item Erstellung einer verständlichen Benutzeranleitung.
\end{itemize}

\section{Entwicklungsumgebung}
\begin{itemize}
    \item \textbf{Programmiersprache:} Java
    \item \textbf{IDE:} IntelliJ IDEA
    \item \textbf{Datenbank:} MariaDB
    \item \textbf{Versionsverwaltung:} Git / GitHub
    \item \textbf{Hardware:} Bookii-Stift, Standard-PC (Windows/Linux), 600-DPI-Drucker.
\end{itemize}

\end{document}
